\documentclass[12pt]{article}
\usepackage[a4paper,margin=1in]{geometry}
\usepackage{iftex}
\ifXeTeX
\usepackage{fontspec}
\defaultfontfeatures{Ligatures=TeX}
\setmainfont{Times New Roman}
\else
\usepackage[T1]{fontenc}
\usepackage[utf8]{inputenc}
\usepackage{newtxtext}
\usepackage{newtxmath}
\fi
\usepackage{enumitem}
\usepackage{hyperref}
\usepackage{parskip}
\usepackage{microtype}
\usepackage[normalem]{ulem}
\usepackage{xcolor}

\newcommand{\practiceid}[1]{\texttt{#1}}
\newlength{\readinggutter}
\setlength{\readinggutter}{2.8em}
\newlength{\readingfirstindent}
\setlength{\readingfirstindent}{1.5em}
\newlength{\questionblockgap}
\setlength{\questionblockgap}{0.45\baselineskip}
\newlength{\exampleblockgap}
\setlength{\exampleblockgap}{0.65\baselineskip}
\newlength{\passageparagraphgap}
\setlength{\passageparagraphgap}{0.45\baselineskip}
\newenvironment{exampleblock}{\par\vspace{0.35\baselineskip}\begingroup\raggedright\setlength{\parskip}{0pt}\setlength{\parindent}{0pt}}{\par\endgroup\vspace{\exampleblockgap}}
\newcommand{\readinglineindent}[2]{%
\par\noindent%
\hangindent=\dimexpr\readinggutter+0.9em\relax%
\hangafter=1%
\makebox[\readinggutter][r]{\scriptsize\color{gray}#1}\hspace{0.9em}\hspace*{\readingfirstindent}#2%
\par}
\newcommand{\readingline}[2]{%
\par\noindent%
\hangindent=\dimexpr\readinggutter+0.9em\relax%
\hangafter=1%
\makebox[\readinggutter][r]{\scriptsize\color{gray}#1}\hspace{0.9em}#2%
\par}

\setlength{\parindent}{0pt}
\setlength{\parskip}{0.5em}
\setlist[enumerate]{topsep=0.25em,itemsep=0.18em,parsep=0pt,partopsep=0pt}
\setlength{\emergencystretch}{2em}

\begin{document}
\section*{TOEFL Practice 07 - Tryout}
Source of truth: DOCX only.
\section*{Listening Comprehension}
\subsection*{Part A}
DIRECTIONS: In Part A you will hear short conversations between two speakers. At the end of each conversation, a third speaker will ask a question about what the first two speakers said. Each conversation and each question will be spoken only one time. Therefore, you must listen carefully to understand what each speaker says. After you hear a conversation and the question, read the four choices and select the one that is the best answer to the question the speaker asked. Then, on your answer. sheet, find the number of the question and blacken the space that corresponds to the letter for the answer you have chosen. Blacken the space completely so that the letter inside the space does not show.\par
\begin{exampleblock}
Listen to the following example.\par
On the recording, you hear:\par
(Man) Does the car need to be filled?\par
(Woman) Mary stopped at the gas station on her way home.\par
(Narrator) What does the woman mean?\par
In your test book, you will read:\par
(A) Mary bought some food.\par
(B) Mary had car trouble.\par
(C) Mary went shopping.\par
(D) Mary bought some gas.\par
From the conversation you learn that Mary stopped at the gas station on her way home. The best answer to the question "Does the car need to be filled?" is (D), "Mary bought some gas." Therefore, the correct answer is (D).\par
Now let us begin Part A with question number 1.\par
\end{exampleblock}
\noindent\textbf{1.}
\begin{enumerate}[label=(\Alph*),leftmargin=*,itemsep=0.2em]
\item He is happy to see the woman again.
\item He is glad the woman went running.
\item He is glad that he can return a favor.
\item He wants a piece of the cake she is baking.
\end{enumerate}
\par
\vspace{0.25\baselineskip}
\noindent\textbf{2.}
\begin{enumerate}[label=(\Alph*),leftmargin=*,itemsep=0.2em]
\item A hair stylist
\item An electrician
\item A die maker
\item A sales clerk
\end{enumerate}
\par
\vspace{0.25\baselineskip}
\noindent\textbf{3.}
\begin{enumerate}[label=(\Alph*),leftmargin=*,itemsep=0.2em]
\item A TV show
\item A newspaper
\item A newscast
\item A schedule
\end{enumerate}
\par
\vspace{0.25\baselineskip}
\noindent\textbf{4.}
\begin{enumerate}[label=(\Alph*),leftmargin=*,itemsep=0.2em]
\item The man is 20 minutes late.
\item The man came to the wrong place again.
\item The man has lost his pen several times.
\item It's too dark to look for the pen.
\end{enumerate}
\par
\vspace{0.25\baselineskip}
\noindent\textbf{5.}
\begin{enumerate}[label=(\Alph*),leftmargin=*,itemsep=0.2em]
\item In Seattle
\item In his office
\item On a plane
\item In San Francisco
\end{enumerate}
\par
\vspace{0.25\baselineskip}
\noindent\textbf{6.}
\begin{enumerate}[label=(\Alph*),leftmargin=*,itemsep=0.2em]
\item Pay for new clothes
\item Exercise regularly
\item Have his clothes altered
\item Weigh himself every day
\end{enumerate}
\par
\vspace{0.25\baselineskip}
\noindent\textbf{7.}
\begin{enumerate}[label=(\Alph*),leftmargin=*,itemsep=0.2em]
\item She is a guest here.
\item She is still thinking.
\item She is finished ordering.
\item She is guessing.
\end{enumerate}
\par
\vspace{0.25\baselineskip}
\noindent\textbf{8.}
\begin{enumerate}[label=(\Alph*),leftmargin=*,itemsep=0.2em]
\item Tracey is very reliable.
\item Tracey doesn't mention Marsha.
\item Tracey is less dependable than Marsha.
\item Tracey is more striking than Marsha.
\end{enumerate}
\par
\vspace{0.25\baselineskip}
\noindent\textbf{9.}
\begin{enumerate}[label=(\Alph*),leftmargin=*,itemsep=0.2em]
\item He is joking.
\item He is not kidding.
\item He agrees with the woman.
\item He likes motorcycle races.
\end{enumerate}
\par
\vspace{0.25\baselineskip}
\noindent\textbf{10.}
\begin{enumerate}[label=(\Alph*),leftmargin=*,itemsep=0.2em]
\item Plan a vacation
\item Buy a camera
\item Take a trip
\item Leave the camera
\end{enumerate}
\par
\vspace{0.25\baselineskip}
\noindent\textbf{11.}
\begin{enumerate}[label=(\Alph*),leftmargin=*,itemsep=0.2em]
\item He is unloading the truck.
\item He has a good mind.
\item He is happy to help the woman.
\item He is also bothered by the noise.
\end{enumerate}
\par
\vspace{0.25\baselineskip}
\noindent\textbf{12.}
\begin{enumerate}[label=(\Alph*),leftmargin=*,itemsep=0.2em]
\item Cancel her appointment
\item Wait until June
\item See the doctor now
\item Make another appointment
\end{enumerate}
\par
\vspace{0.25\baselineskip}
\noindent\textbf{13.}
\begin{enumerate}[label=(\Alph*),leftmargin=*,itemsep=0.2em]
\item She doesn't like it when people look at her.
\item She doesn't intend to lend the man the money.
\item She works in the loan department at the bank.
\item Unlike the man, she has paid her tuition already.
\end{enumerate}
\par
\vspace{0.25\baselineskip}
\noindent\textbf{14.}
\begin{enumerate}[label=(\Alph*),leftmargin=*,itemsep=0.2em]
\item The cut should be made differently than was proposed.
\item There is ample time to analyze the data.
\item The time for data analysis may not be sufficient.
\item The market survey won't be completed in time.
\end{enumerate}
\par
\vspace{0.25\baselineskip}
\noindent\textbf{15.}
\begin{enumerate}[label=(\Alph*),leftmargin=*,itemsep=0.2em]
\item She knows about the road construction.
\item She can't find an alternate route home.
\item She wants more information about the traffic.
\item She wants the man to tell her about his problems.
\end{enumerate}
\par
\vspace{0.25\baselineskip}
\noindent\textbf{16.}
\begin{enumerate}[label=(\Alph*),leftmargin=*,itemsep=0.2em]
\item They should. have stayed in their hotel.
\item They should have checked their route.
\item The stadium is difficult to find.
\item The map was lost while they were driving.
\end{enumerate}
\par
\vspace{0.25\baselineskip}
\noindent\textbf{17.}
\begin{enumerate}[label=(\Alph*),leftmargin=*,itemsep=0.2em]
\item Discuss buying the house
\item Talk to the buyer
\item Appraise the property
\item Play volleyball
\end{enumerate}
\par
\vspace{0.25\baselineskip}
\noindent\textbf{18.}
\begin{enumerate}[label=(\Alph*),leftmargin=*,itemsep=0.2em]
\item Looking for something else for the children to do
\item Finding other means of transportation
\item Canceling the trip to the newspaper
\item Taking the trip on another day of the week
\end{enumerate}
\par
\vspace{0.25\baselineskip}
\noindent\textbf{19.}
\begin{enumerate}[label=(\Alph*),leftmargin=*,itemsep=0.2em]
\item Playing cards
\item Writing letters
\item Playing a board game
\item Memorizing vocabulary
\end{enumerate}
\par
\vspace{0.25\baselineskip}
\noindent\textbf{20.}
\begin{enumerate}[label=(\Alph*),leftmargin=*,itemsep=0.2em]
\item She is surprised.
\item She is excited.
\item She is angry.
\item She is disgusted.
\end{enumerate}
\par
\vspace{0.25\baselineskip}
\noindent\textbf{21.}
\begin{enumerate}[label=(\Alph*),leftmargin=*,itemsep=0.2em]
\item He cannot eat dinner alone.
\item He came in with two other people:
\item Three people will have dinner.
\item A table for two people is needed.
\end{enumerate}
\par
\vspace{0.25\baselineskip}
\noindent\textbf{22.}
\begin{enumerate}[label=(\Alph*),leftmargin=*,itemsep=0.2em]
\item The salesman was not helpful.
\item The salesman was excellent.
\item The man didn't need help.
\item The service was a whole lot better.
\end{enumerate}
\par
\vspace{0.25\baselineskip}
\noindent\textbf{23.}
\begin{enumerate}[label=(\Alph*),leftmargin=*,itemsep=0.2em]
\item The ball is in the front hall.
\item The light doesn't work.
\item She can smell something burning.
\item She cannot find the light.
\end{enumerate}
\par
\vspace{0.25\baselineskip}
\noindent\textbf{24.}
\begin{enumerate}[label=(\Alph*),leftmargin=*,itemsep=0.2em]
\item He is tired of sandwiches.
\item He's already had lunch.
\item He's gained weight.
\item It's too late for sandwiches.
\end{enumerate}
\par
\vspace{0.25\baselineskip}
\noindent\textbf{25.}
\begin{enumerate}[label=(\Alph*),leftmargin=*,itemsep=0.2em]
\item She couldn't hear the band.
\item The band was not very good
\item She plays better than the band
\item She liked it very much.
\end{enumerate}
\par
\vspace{0.25\baselineskip}
\noindent\textbf{26.}
\begin{enumerate}[label=(\Alph*),leftmargin=*,itemsep=0.2em]
\item He doesn't want to swim.
\item The pool is not in the city.
\item He is a city resident.
\item Nine years is a long time.
\end{enumerate}
\par
\vspace{0.25\baselineskip}
\noindent\textbf{27.}
\begin{enumerate}[label=(\Alph*),leftmargin=*,itemsep=0.2em]
\item They have an hour to spare.
\item The airport is not that far.
\item The clock should be repaired.
\item The clock is not accurate.
\end{enumerate}
\par
\vspace{0.25\baselineskip}
\noindent\textbf{28.}
\begin{enumerate}[label=(\Alph*),leftmargin=*,itemsep=0.2em]
\item The woman is a good guitar player.
\item The woman keeps changing her mind about hobbies.
\item He can teach the woman to take pictures and play guitar.
\item Photography is similar to playing guitar.
\end{enumerate}
\par
\vspace{0.25\baselineskip}
\noindent\textbf{29.}
\begin{enumerate}[label=(\Alph*),leftmargin=*,itemsep=0.2em]
\item She will drive the man.
\item Her car doesn't work.
\item The man can borrow her car.
\item She plans to go shopping
\end{enumerate}
\par
\vspace{0.25\baselineskip}
\noindent\textbf{30.}
\begin{enumerate}[label=(\Alph*),leftmargin=*,itemsep=0.2em]
\item The instructor has a responsibility to help the man.
\item The assignments are unreasonably difficult.
\item The course is too advanced for the man.
\item Students are expected to do the assignments.
\end{enumerate}
\par
\vspace{0.25\baselineskip}
\subsection*{Part B}
DIRECTIONS: In this part of the test, you will hear longer conversations. After each conversation, you will hear several questions. The conversations and questions will not be repeated. After you hear a question, read the four possible answers in. your test book and select. The best answer. Then, on your answer sheet, find the number of the question and fill in the space that corresponds to the letter of the answer you have chosen. Remember, you are not allowed to take notes or write in your test book.\par
\begin{exampleblock}
Listen to the following example:\par
You will hear:\par
You will read:\par
(A) He has changed jobs.\par
(B) He has two children.\par
(C) He has two jobs.\par
(D) He is looking for a job.\par
From the conversation you learn that Tom has taken an additional job. The best answer to the question "Why is Tom tired?" is (C), "He has two jobs." Therefore the correct answer is (C).\par
\end{exampleblock}
\noindent\textbf{31.}
\begin{enumerate}[label=(\Alph*),leftmargin=*,itemsep=0.2em]
\item Gathering information about colleges
\item Choosing a good graduate school
\item Writing letters to apply for a job
\item Discovering essential facts
\end{enumerate}
\par
\vspace{0.25\baselineskip}
\noindent\textbf{32.}
\begin{enumerate}[label=(\Alph*),leftmargin=*,itemsep=0.2em]
\item Tuition and fees
\item Motivation
\item Job experience
\item Computers
\end{enumerate}
\par
\vspace{0.25\baselineskip}
\noindent\textbf{33.}
\begin{enumerate}[label=(\Alph*),leftmargin=*,itemsep=0.2em]
\item Going to college
\item Taking classes
\item Asking questions
\item Writing brochures
\end{enumerate}
\par
\vspace{0.25\baselineskip}
\noindent\textbf{34.}
\begin{enumerate}[label=(\Alph*),leftmargin=*,itemsep=0.2em]
\item Studying is similar to working.
\item They learn to work with computers.
\item In small classes, they prepare for jobs.
\item They take jobs available on campus.
\end{enumerate}
\par
\vspace{0.25\baselineskip}
\noindent\textbf{35.}
\begin{enumerate}[label=(\Alph*),leftmargin=*,itemsep=0.2em]
\item Mostly for herself
\item Primarily for her children
\item To guard the house
\item To take it to Scotland
\end{enumerate}
\par
\vspace{0.25\baselineskip}
\noindent\textbf{36.}
\begin{enumerate}[label=(\Alph*),leftmargin=*,itemsep=0.2em]
\item They are disciplined and reliable.
\item They cannot stay alone at home.
\item Children can take them to school,
\item Fifteen collies can herd cattle.
\end{enumerate}
\par
\vspace{0.25\baselineskip}
\noindent\textbf{37.}
\begin{enumerate}[label=(\Alph*),leftmargin=*,itemsep=0.2em]
\item Golden retrievers
\item Cocker spaniels
\item Terriers
\item Foxhounds
\end{enumerate}
\par
\vspace{0.25\baselineskip}
\noindent\textbf{38.}
\begin{enumerate}[label=(\Alph*),leftmargin=*,itemsep=0.2em]
\item Terriers have thick brown coats.
\item The family has little time to spend with it.
\item She should take of her children, instead of the dog.
\item Her children didnt like terriers as pets.
\end{enumerate}
\par
\vspace{0.25\baselineskip}
\subsection*{Part C}
DIRECTIONS: In Part C you will hear short lectures and extended conversations. At the end of each, you will be asked several questions. Each lecture or conversation and each question will be spoken only one time. For this reason, you must listen carefully to understand what each speaker says. After you hear a question, read the four choices and select the one that best answers the question the speaker asked. Then, on your answer sheet, find the number of the question and blacken the space that contains the letter for the answer you have chosen. Answer all questions according to what is stated or implied in the lecture or conversation.\par
\begin{exampleblock}
Listen to this sample talk.\par
You will hear:\par
Now listen, to the following example.\par
You will hear:\par
You will read:\par
(A) By cars and carriages\par
(B) By bicycles, trains, and carriages\par
(C) On foot and by boat\par
(D) On board ships and trains\par
The best answer to the question "According to the speaker, how did people travel before the invention of the automobile?" is (B), "By bicycles, trains, and carriages." Therefore, the correct answer is (B). You are not allowed to make notes during the test.\par
\end{exampleblock}
\noindent\textbf{39.}
\begin{enumerate}[label=(\Alph*),leftmargin=*,itemsep=0.2em]
\item Video rental business
\item Large motion picture studios
\item Vehicle rental agencies
\item Convenience and department stores
\end{enumerate}
\par
\vspace{0.25\baselineskip}
\noindent\textbf{40.}
\begin{enumerate}[label=(\Alph*),leftmargin=*,itemsep=0.2em]
\item Between 300 and 600
\item Between 5,000 and 7,000
\item Approximately 60,000
\item Over 90,000
\end{enumerate}
\par
\vspace{0.25\baselineskip}
\noindent\textbf{41.}
\begin{enumerate}[label=(\Alph*),leftmargin=*,itemsep=0.2em]
\item Rental space can be hard to find.
\item It's competitive and volatile.
\item Few stores in big cities rent videos.
\item It includes fast-food restaurants.
\end{enumerate}
\par
\vspace{0.25\baselineskip}
\noindent\textbf{42.}
\begin{enumerate}[label=(\Alph*),leftmargin=*,itemsep=0.2em]
\item Motion picture studios
\item Distribution networks
\item Convenience stores
\item Big chains of stores
\end{enumerate}
\par
\vspace{0.25\baselineskip}
\noindent\textbf{43.}
\begin{enumerate}[label=(\Alph*),leftmargin=*,itemsep=0.2em]
\item It extends the existing supply of minerals.
\item t contributes to pollution and waste.
\item Iron and aluminum are in short supply.
\item Mining aluminum costs a great deal.
\end{enumerate}
\par
\vspace{0.25\baselineskip}
\noindent\textbf{44.}
\begin{enumerate}[label=(\Alph*),leftmargin=*,itemsep=0.2em]
\item 40 percent
\item 65 percent
\item One quarter
\item One half
\end{enumerate}
\par
\vspace{0.25\baselineskip}
\noindent\textbf{45.}
\begin{enumerate}[label=(\Alph*),leftmargin=*,itemsep=0.2em]
\item They have to pay taxes to recycle.
\item They are used to a throwaway lifestyle.
\item Environment is not a serious concern in the U.S.
\item Recycling does not provide many benefits.
\end{enumerate}
\par
\vspace{0.25\baselineskip}
\noindent\textbf{46.}
\begin{enumerate}[label=(\Alph*),leftmargin=*,itemsep=0.2em]
\item They encourage it.
\item They don't like it.
\item They make money from it.
\item They discuss it.
\end{enumerate}
\par
\vspace{0.25\baselineskip}
\noindent\textbf{47.}
\begin{enumerate}[label=(\Alph*),leftmargin=*,itemsep=0.2em]
\item A description of roses
\item The design of a colorful rose garden
\item Growing roses in various climates
\item Caring for shrubs and flowers
\end{enumerate}
\par
\vspace{0.25\baselineskip}
\noindent\textbf{48.}
\begin{enumerate}[label=(\Alph*),leftmargin=*,itemsep=0.2em]
\item The higher the rating, the fewer climates are suitable for a rose.
\item The higher the rating, the more areas are appropriate for growing a rose.
\item The highest rated rose is the most difficult to grow.
\item The highest rank is given to specially cultivated roses.
\end{enumerate}
\par
\vspace{0.25\baselineskip}
\noindent\textbf{49.}
\begin{enumerate}[label=(\Alph*),leftmargin=*,itemsep=0.2em]
\item Their flowers fail to open.
\item Their petals are limited in number and color.
\item Their flower production drops in summer heat.
\item They need to be watered frequently.
\end{enumerate}
\par
\vspace{0.25\baselineskip}
\noindent\textbf{50.}
\begin{enumerate}[label=(\Alph*),leftmargin=*,itemsep=0.2em]
\item Develop hardy varieties of roses.
\item Replace roses every spring.
\item Take a trip to a municipal garden.
\item Protect the shrubs from extreme cold.
\end{enumerate}
\par
\vspace{0.25\baselineskip}
\section*{Structure and Written Expression}
\subsection*{Sentence Completion}
This section is designed to test your ability to recognize language structures that are appropriate in standard written English. The questions in this section belong to two types, each of which has special directions.\par
DIRECTIONS: Questions 1-15 are partial sentences. Below each sentence you will see four words or phrases, marked (A), (B), (C), and (D). Select the one word or phrase that best completes the sentence. Then, on your answer sheet, find the number of the question and fill in the space that contains the letter for the answer you have chosen. Fill in the space completely.\par
\begin{exampleblock}
\textbf{EXAMPLE I}\par
\noindent Drying flowers is the best way......them.
\begin{enumerate}[label=(\Alph*),leftmargin=*,itemsep=0.2em]
\item to preserve
\item by preserving
\item preserve
\item preserved
\end{enumerate}
\vspace{0.25\baselineskip}
The sentence should state, "Drying flowers is the best way to preserve them." Therefore, the correct answer is (A).\par
\end{exampleblock}
\begin{exampleblock}
\textbf{EXAMPLE II}\par
\noindent Many American universities......as small, private colleges.
\begin{enumerate}[label=(\Alph*),leftmargin=*,itemsep=0.2em]
\item begun
\item beginning
\item began
\item for the beginning
\end{enumerate}
\vspace{0.25\baselineskip}
The sentence should state, "Many American universities began as small, private colleges." Therefore, the correct answer is (C). After you read the directions, begin work on the questions.\par
\end{exampleblock}
\noindent\textbf{1.} Beginning in the 1960s, national studies have surveyed changes in the values of college students to identify......in college education.
\begin{enumerate}[label=(\Alph*),leftmargin=*,itemsep=0.2em]
\item that it is important to them
\item what is important to them
\item to them what is the importance
\item what importance to them
\end{enumerate}
\par
\vspace{0.25\baselineskip}
\noindent\textbf{2.} Bacteria that......in ocean waters transport other organisms to new habitable environments and sources of nourishment.
\begin{enumerate}[label=(\Alph*),leftmargin=*,itemsep=0.2em]
\item free travel
\item traveling freely
\item travel freely
\item freely travels
\end{enumerate}
\par
\vspace{0.25\baselineskip}
\noindent\textbf{3.} Two objects, both with positive or both with negative charges, are repelled. while objects with opposite charges......
\begin{enumerate}[label=(\Alph*),leftmargin=*,itemsep=0.2em]
\item are attracting the other
\item attracted to each other
\item to each other attract them
\item are attracted to each other
\end{enumerate}
\par
\vspace{0.25\baselineskip}
\noindent\textbf{4.} Roman law, one of......with a continuous history, spread to many countries during the Roman conquest of Europe and Asia.
\begin{enumerate}[label=(\Alph*),leftmargin=*,itemsep=0.2em]
\item the oldest existing systems
\item systems, the oldest existing
\item the existence, the oldest system
\item existing the oldest system
\end{enumerate}
\par
\vspace{0.25\baselineskip}
\noindent\textbf{5.} Although the eucalyptus is native to Australia,......in the U.S. belongs to the eucalyptus family, which comprises over 500 varieties.
\begin{enumerate}[label=(\Alph*),leftmargin=*,itemsep=0.2em]
\item a large quantity of timber is harvested
\item largely, a quantity of timber is harvested
\item a large quantity of timber harvested
\item a quantity of timber is large, harvested
\end{enumerate}
\par
\vspace{0.25\baselineskip}
\noindent\textbf{6.} Florida, mainly a peninsula jutting into the Atlantic Ocean,......during the economic expansion of the 1990s.
\begin{enumerate}[label=(\Alph*),leftmargin=*,itemsep=0.2em]
\item has been experienced, its fastest growth
\item has experienced, it grows the fastest
\item experienced its fastest growth
\item its fastest growth has been experiencing
\end{enumerate}
\par
\vspace{0.25\baselineskip}
\noindent\textbf{7.} When an intersecting plane......of a conical surface, the resulting curve is called the hyperbola.
\begin{enumerate}[label=(\Alph*),leftmargin=*,itemsep=0.2em]
\item cuts both sides
\item cutting both side
\item both sides are cut
\item cut in both sides
\end{enumerate}
\par
\vspace{0.25\baselineskip}
\noindent\textbf{8.} After......California, David Hockney contributed to the Pop art trend and developed his figurative style in painting portraits.
\begin{enumerate}[label=(\Alph*),leftmargin=*,itemsep=0.2em]
\item he, moving to
\item to moving
\item moved to
\item moving to
\end{enumerate}
\par
\vspace{0.25\baselineskip}
\noindent\textbf{9.} When the demand for merchandise exceeds capacity, customers......, if no backup inventory is available.
\begin{enumerate}[label=(\Alph*),leftmargin=*,itemsep=0.2em]
\item must turn away
\item must turn them away
\item must be turned away
\item turn away and must be
\end{enumerate}
\par
\vspace{0.25\baselineskip}
\noindent\textbf{10.} Although......to a public office, Robert Moses served as the New York secretary of state between 1927 and 1928 and led the state redevelopment program.
\begin{enumerate}[label=(\Alph*),leftmargin=*,itemsep=0.2em]
\item he never elected
\item electing never
\item never electing
\item never elected
\end{enumerate}
\par
\vspace{0.25\baselineskip}
\noindent\textbf{11.} Ostriches are flightless birds that can be......2.5 meters with exceptionally strong legs.
\begin{enumerate}[label=(\Alph*),leftmargin=*,itemsep=0.2em]
\item as tall
\item so tall as
\item tall so as
\item as tall as
\end{enumerate}
\par
\vspace{0.25\baselineskip}
\noindent\textbf{12.} Sandy beaches are popular recreational locations where people come......
\begin{enumerate}[label=(\Alph*),leftmargin=*,itemsep=0.2em]
\item to swim, play sports, and sunbathe
\item to swimming, playing sports, and sunbathing
\item for a swim, playing sports and sunbathing
\item swim, play sports, and sunbathe
\end{enumerate}
\par
\vspace{0.25\baselineskip}
\noindent\textbf{13.} Ballet dances appear to perform many movements......
\begin{enumerate}[label=(\Alph*),leftmargin=*,itemsep=0.2em]
\item that is unnatural in the body
\item that are unnatural for the body
\item where the body is unnatural
\item when the body is being unnatural
\end{enumerate}
\par
\vspace{0.25\baselineskip}
\noindent\textbf{14.} Export credit is a financial guarantee usually provided by a bank to enable companies......when the payment is delayed or risky.
\begin{enumerate}[label=(\Alph*),leftmargin=*,itemsep=0.2em]
\item for the goods exported
\item to export goods
\item to exporting goods
\item exporting the goods to
\end{enumerate}
\par
\vspace{0.25\baselineskip}
\noindent\textbf{15.} Nightingales that can be heard at night as well as by day migrate to Africa in winters and Europe in summers......in abundant supply.
\begin{enumerate}[label=(\Alph*),leftmargin=*,itemsep=0.2em]
\item when is food
\item when food is
\item their food is
\item there food is
\end{enumerate}
\par
\vspace{0.25\baselineskip}
\subsection*{Error Identification}
DIRECTIONS: In questions 16-40 every sentence has four words or phrases that are underlined. The four underlined portions of each sentence are marked (A), (B), (C), and (D). Identify the one word or phrase that makes the sentence incorrect. Then, on your answer sheet, find the number of the question and fill in the space that contains the letter for the answer you have chosen. Fill in the space completely.\par
\begin{exampleblock}
\textbf{EXAMPLE I}\par
Christopher Columbus \uline{has sailed}(A) from Europe in 1492 and \uline{discovered a}(B) \uline{new land}(C) he \uline{thought to}(D) be India.\par
The sentence should state, "Christopher Columbus sailed from Europe in 1492 and discovered a new land he thought to be India." Therefore, you should choose answer (A).\par
\end{exampleblock}
\begin{exampleblock}
\textbf{EXAMPLE II}\par
\uline{As the roles}(A) of people in society change, \uline{so does}(B) the \uline{rules of}(C) conduct in certain \uline{situations}(D).\par
The sentence should state, "As the roles of people in society change, so do the rules of\par
conduct in certain situations." Therefore, you should choose answer (B).\par
\end{exampleblock}
\noindent\textbf{16.} \uline{Having completing}(A)his education in Harvard, Wendell Phillips became \uline{an outspoken}(B)abolitionist and a \uline{widely-traveled}(C)lecturer, who also \uline{supported labor}(D)unionization.
\par
\vspace{\questionblockgap}
\noindent\textbf{17.} Plant hormones \uline{have shown}(A)to produce a remarkable \uline{effect on}(B)\uline{plant growth}(C), flowering, and \uline{the generation of}(D)fruit.
\par
\vspace{\questionblockgap}
\noindent\textbf{18.} Lately, motorcycle \uline{race courses}(A)\uline{incorporated}(B)oval-shaped dirt tracks, \uline{hill climbs}(C), and natural terrains with \uline{added artificial}(D)hazards.
\par
\vspace{\questionblockgap}
\noindent\textbf{19.} Passports \uline{are legal}(A)documents \uline{issued by}(B)governments to give their bearers an authorization to \uline{cross nation}(C)borders and \uline{seek protection}(D)in an emergency.
\par
\vspace{\questionblockgap}
\noindent\textbf{20.} \uline{Body of still}(A)water in ground depressions are common \uline{in formerly}(B)glaciated areas, near \uline{slow-flowing}(C)rivers, and \uline{in low lands}(D)near sea shores.
\par
\vspace{\questionblockgap}
\noindent\textbf{21.} It is \uline{not known}(A)why \uline{many seeds undergo}(B)a period of dormancy \uline{even during the times when}(C)conditions for \uline{their growth is}(D)favorable.
\par
\vspace{\questionblockgap}
\noindent\textbf{22.} \uline{The lowest layer}(A)of the atmosphere, \uline{heated by}(B)the Earth, \uline{consisting of}(C)bound atoms of nitrogen and oxygen, \uline{with a small}(D)quantity of argon.
\par
\vspace{\questionblockgap}
\noindent\textbf{23.} Alvin Ailey's choreography \uline{reflected}(A)his southern roots and \uline{center}(B)on \uline{fragments of}(C)folk songs and jazz that was \uline{innovative}(D)in his time.
\par
\vspace{\questionblockgap}
\noindent\textbf{24.} The \uline{wire-wound}(A)coil \uline{in a simple}(B)generator \uline{rotates between}(C)the \uline{pole of}(D)a permanent magnet.
\par
\vspace{\questionblockgap}
\noindent\textbf{25.} Milwaukee, \uline{settled in}(A)1818, \uline{drawn a}(B)large \uline{number of}(C)immigrants who \uline{arrived in}(D)search of jobs and inexpensive housing.
\par
\vspace{\questionblockgap}
\noindent\textbf{26.} J.R.R. Tolkien's book for children, titled The Hobbit, \uline{describes the adventures}(A)of Bilbo \uline{in a fantastic world}(B), populated by dragons, dwarfs, and \uline{another mythical}(C)creatures, \uline{including}(D)a wizard.
\par
\vspace{\questionblockgap}
\noindent\textbf{27.} Studies of Mars \uline{indicate that}(A)enough \uline{water might}(B)be collected on \uline{the planet's}(C)surface \uline{sustain prolonged}(D)missions by human space crews.
\par
\vspace{\questionblockgap}
\noindent\textbf{28.} North Dakota, \uline{rich in}(A)natural gas and lignite deposits, \uline{remains}(B)\uline{the most rural}(C)of all states, with 90\% \uline{of land}(D)under agricultural cultivation.
\par
\vspace{\questionblockgap}
\noindent\textbf{29.} Geological dating can be carried out \uline{by measuring}(A)how much radioactive \uline{elements of a}(B)rock \uline{have change}(C)since the rock \uline{was formed}(D).
\par
\vspace{\questionblockgap}
\noindent\textbf{30.} \uline{Parallel beams}(A)of light converge and \uline{can be}(B)eventually brought into focus \uline{when pass}(C)through a convex lens \uline{to produce}(D)a real image on a screen.
\par
\vspace{\questionblockgap}
\noindent\textbf{31.} If the air \uline{temperature falls}(A)below the \uline{freezing point}(B), the dew \uline{will freeze}(C), and the water vapor \uline{can condenses}(D)directly into ice.
\par
\vspace{\questionblockgap}
\noindent\textbf{32.} \uline{Although different}(A)in size and shape, most ferns are perennial and \uline{spreading}(B)by \uline{low-growing}(C)root systems and spores \uline{transported}(D)by wind.
\par
\vspace{\questionblockgap}
\noindent\textbf{33.} Walt Disney and his brother established \uline{their animation}(A)studio in Hollywood in 1923, and the \uline{first animation}(B)cartoon with Mickey Mouse \uline{appeared in}(C)\uline{black and white}(D)in 1928.
\par
\vspace{\questionblockgap}
\noindent\textbf{34.} For centuries, philosophers and \uline{artists alike}(A)\uline{have debated}(B)the meaning of beauty and questioned \uline{whether it real}(C)or \uline{only perceived}(D).
\par
\vspace{\questionblockgap}
\noindent\textbf{35.} In some musical compositions, regular \uline{alternations}(A)of the beat \uline{are avoided}(B)without repetition \uline{so that}(C)the listener will not become \uline{accustomed by}(D)them.
\par
\vspace{\questionblockgap}
\noindent\textbf{36.} Margaret Mitchell, the author of Gone with the Wind, \uline{had researched}(A)the Civil War and then \uline{had written}(B)\uline{the story about}(C)Georgia and the period \uline{that followed}(D)the war.
\par
\vspace{\questionblockgap}
\noindent\textbf{37.} \uline{As the sign}(A)languages of \uline{the deaf}(B)show, \uline{the sounds}(C)of speech are not required \uline{to expressing}(D)meaningful structures.
\par
\vspace{\questionblockgap}
\noindent\textbf{38.} In 1831, \uline{English physicist}(A), Michael Faraday, discovered \uline{the principle of}(B)electromagnetic induction, \uline{which the}(C)electric motor and the electric generator \uline{are based}(D).
\par
\vspace{\questionblockgap}
\noindent\textbf{39.} \uline{Aspiring to}(A)eliminate slavery \uline{entirely}(B), Francois Toussaint organized an army \uline{consisting mostly of}(C)illiterate slaves into \uline{a power}(D)military force.
\par
\vspace{\questionblockgap}
\noindent\textbf{40.} \uline{The human}(A)rib cage protects the lungs and \uline{the heart allows}(B)the chest \uline{to expand}(C)and \uline{contract easily}(D).
\par
\vspace{\questionblockgap}
\section*{Reading Comprehension}
\begin{center}
\textbf{Time 55 minutes}
\end{center}
DIRECTIONS: In this section you will read several passages. Each passage is followed by a series of questions. For questions 1-50, you need to select the best answer, (A), (B), (C), or (D), to each question. Then, on your answer sheet, find the number of the question and fill in the space that contains the letter of the answer you have selected. Fill in the space completely. Answer all questions following a passage on the basis of what is stated or implied in the passage.\par
\begin{exampleblock}
\small
\sloppy
Read the following passage:\par
A tomahawk is a small ax used as a tool and a weapon by the North American Indian\par
tribes. An average tomahawk was not very long and did not weigh a great deal. Originally,\par
the head of the tomahawk was made of a shaped stone or an animal bone and was mounted on\par
Line a wooden handle. After the arrival of the European settlers, the Indians began to use toma-\par
hawks with iron heads. Indian males and females of all ages used tomahawks to chop and cut wood, pound stakes into the ground to put up wigwams, and do many other chores. Indian warriors relied on tomahawks as weapons and even threw them at their enemies. Some types of tomahawks were used in religious ceremonies. Contemporary American idioms reflect\par
this aspect of American heritage.\par
\end{exampleblock}
\begin{exampleblock}
\textbf{EXAMPLE I}\par
\noindent Early tomahawk heads were made of
\begin{enumerate}[label=(\Alph*),leftmargin=*,itemsep=0.2em]
\item stone or bone
\item wood or sticks
\item European iron
\item religious weapons
\end{enumerate}
\vspace{0.25\baselineskip}
According to the passage, early tomahawk heads were made of stone or bone. Therefore, the correct answer is (A).\par
\end{exampleblock}
\begin{exampleblock}
\textbf{EXAMPLE II}\par
\noindent How has the Indian use of tomahawks affected American daily life today?
\begin{enumerate}[label=(\Alph*),leftmargin=*,itemsep=0.2em]
\item Tomahawks are still used as weapons.
\item Tomahawks are used as tools for.certain jobs.
\item Contemporary language refers to tomahawks.
\item Indian tribes cherish tomahawks as heirlooms.
\end{enumerate}
\vspace{0.25\baselineskip}
The passage states, "Contemporary American idioms reflect this aspect of American heritage." The correct answer is (C). After you read the directions, begin work on the questions.\par
\end{exampleblock}
\subsection*{Questions 1-12}
\begingroup
\small
\sloppy
\setlength{\parskip}{0pt}
\setlength{\parindent}{0pt}
In the course of history, human inventions have dramatically increased the average\par
amount of energy available for use per person. Primitive peoples in cold regions burned\par
wood and animal dung to heat their caves, cook food, and drive off animals by fire. The first\par
step toward the developing of more efficient fuels was taken when people discovered that\par
they could use vegetable oils and animal fats in lieu of gathered or cut wood. Charcoal gave\par
off a more intensive heat than wood and was more easily obtainable than organic fats. The\par
Greeks first began to use coal for metal smelting in the 4th century, but it did not come into\par
extensive use until the Industrial Revolution.\par
In the 1700s, at the beginning of the Industrial Revolution, most energy used in the\par
United States and other nations undergoing industrialization was obtained from perpetual\par
and renewable sources, such as wood, water streams, domesticated animal labor, and wind.\par
These were. predominantly locally available supplies. By mid-1800s, 91 percent of all com-\par
mercial energy consumed in the United States and European countries was obtained from\par
wood. However, at the beginning of the 20th century, coal became a major energy source\par
and replaced wood in industrializing countries. Although in most regions and climate.\par
zones wood was more readily accessible than coal, the latter represents a more concentrated\par
source of energy. In 1910, natural gas and oil firmly replaced coal as the main source of fuel\par
because they are lighter and, therefore, cheaper to transport. They burned more cleanly than\par
coal and polluted less. Unlike coal, oil could be refined to manufacture liquid fuels for\par
vehicles, a very important consideration in the early 1900s, when the automobile arrived on\par
the scene.\par
By 1984, nonrenewable fossil fuels, such as oil, coal; and natural gas, provided over 82\par
percent of the commercial and industrial energy used in the world. Small amounts of energy\par
were derived from nuclear fission, and the remaining 16 percent came from burning direct\par
perpetual and renewable fuels, such as biomass. Between 1700 and 1986, a large number of\par
countries shifted from the use of energy from local sources to a centralized generation of hy-\par
dropower and solar energy converted to electricity. The energy derived from nonrenewable\par
fossil fuels has been increasingly produced in one location and transported to another, as is\par
the case with most automobile fuels. In countries with private, rather than public transpor-\par
tation, the age of nonrenewable fuels has created a dependency on a finite resource that will\par
have to be replaced.\par
Alternative fuel sources are numerous, and shale oil and hydrocarbons are just two exam-\par
ples. The extraction of shale oil from large deposits in Asian and European regions has\par
proven to be labor consuming and costly. The resulting product is sulfur- and nitrogen-\par
rich, and large-scale extractions are presently prohibitive. Similarly, the extraction of hydro-\par
carbons from tar sands in Alberta and. Utah is complex. Semi-solid hydrocarbons cannot be\par
easily separated from the sandstone and limestone that carry them, and modern technology\par
is not sufficiently versatile for a large-scale removal of the material. However, both sources\par
of fuel may eventually be needed as petroleum prices continue to rise and limitations in\par
fossil fuel availability make alternative deposits more attractive,\par
\endgroup
\medskip
\noindent\textbf{1.} What is the main topic of the passage?
\begin{enumerate}[label=(\Alph*),leftmargin=*,itemsep=0.2em]
\item Applications of various fuels
\item Natural resources and fossil fuels
\item A history of energy use
\item A historical overview of energy rates
\end{enumerate}
\par
\vspace{0.25\baselineskip}
\noindent\textbf{2.} In line 2, the phrase "per person" is closest in meaning to
\begin{enumerate}[label=(\Alph*),leftmargin=*,itemsep=0.2em]
\item per capita
\item per year
\item per family
\item per day
\end{enumerate}
\par
\vspace{0.25\baselineskip}
\noindent\textbf{3.} It can be inferred from the first paragraph that
\begin{enumerate}[label=(\Alph*),leftmargin=*,itemsep=0.2em]
\item coal mining was essential for primitive peoples
\item the Greeks used coal in industrial production
\item the development of efficient fuels was a gradual process
\item the discovery of efficient fuels was mostly accidental
\end{enumerate}
\par
\vspace{0.25\baselineskip}
\noindent\textbf{4.} In line 5, the phrase "in lieu" is closest in meaning to
\begin{enumerate}[label=(\Alph*),leftmargin=*,itemsep=0.2em]
\item in spite
\item in place
\item in every way
\item in charge
\end{enumerate}
\par
\vspace{0.25\baselineskip}
\noindent\textbf{5.} The author of the passage implies that in the 1700s, sources of energy were
\begin{enumerate}[label=(\Alph*),leftmargin=*,itemsep=0.2em]
\item used for commercial purposes
\item used in various combinations
\item not derived from mineral deposits
\item not always easy to locate
\end{enumerate}
\par
\vspace{0.25\baselineskip}
\noindent\textbf{6.} In line 16, the phrase "the latter" refers to
\begin{enumerate}[label=(\Alph*),leftmargin=*,itemsep=0.2em]
\item wood
\item coal
\item most regions
\item climate zones
\end{enumerate}
\par
\vspace{0.25\baselineskip}
\noindent\textbf{7.} In line 18, the word "They" refers to
\begin{enumerate}[label=(\Alph*),leftmargin=*,itemsep=0.2em]
\item coal and wood
\item main sources of fuel
\item natural gas and oil
\item industrializing countries
\end{enumerate}
\par
\vspace{0.25\baselineskip}
\noindent\textbf{8.} According to the passage, what was the greatest advantage of oil as fuel?
\begin{enumerate}[label=(\Alph*),leftmargin=*,itemsep=0.2em]
\item It was a concentrated source of energy.
\item It was lighter and cheaper than coal.
\item It replaced wood and coal and reduced pollution.
\item It could be converted to automobile fuel.
\end{enumerate}
\par
\vspace{0.25\baselineskip}
\noindent\textbf{9.} According to the passage, the sources of fossil fuels will have to be replaced because
\begin{enumerate}[label=(\Alph*),leftmargin=*,itemsep=0.2em]
\item they need to be transported
\item they are not efficient
\item their use is centralized.
\item their supply is limited
\end{enumerate}
\par
\vspace{0.25\baselineskip}
\noindent\textbf{10.} It can be inferred from the passage that in the early 20th century, energy was obtained primarily from
\begin{enumerate}[label=(\Alph*),leftmargin=*,itemsep=0.2em]
\item fossil fuels
\item nuclear fission
\item hydraulic and solar sources
\item burning biomass
\end{enumerate}
\par
\vspace{0.25\baselineskip}
\noindent\textbf{11.} The author of the passage implies that alternative sources of fuel are currently
\begin{enumerate}[label=(\Alph*),leftmargin=*,itemsep=0.2em]
\item being explored
\item being used for consumption
\item available in few locations
\item examined on a large scale
\end{enumerate}
\par
\vspace{0.25\baselineskip}
\noindent\textbf{12.} In line 35, the word "prohibitive" is closest in meaning to
\begin{enumerate}[label=(\Alph*),leftmargin=*,itemsep=0.2em]
\item prohibited
\item provided
\item too expensive
\item too expedient
\end{enumerate}
\par
\vspace{0.25\baselineskip}
\subsection*{Questions 13-23}
\begingroup
\small
\sloppy
\setlength{\parskip}{0pt}
\setlength{\parindent}{0pt}
Most forms of property are concrete and tangible, such as houses, cars, furniture, or any-\par
thing else that is included in one's possessions. Other forms of property can be intangible,\par
and copyright deals with intangible forms of property. Copyright is a legal protection ex-\par
tended to authors of creative works, for example, books, magazine articles, maps, films,\par
plays, television shows, software, paintings, photographs, music, choreography in dance,\par
and all other forms of intellectual or artistic property.\par
Although the purpose of artistic property is usually public use and enjoyment, copyright\par
establishes the ownership of the creator. When a person buys a copyrighted magazine, it be-\par
longs to this individual as a tangible object. However, the authors of the magazine articles\par
own the research and the writing that went into creating the articles. The right to make and\par
sell or give away copies of books or articles belongs to the authors, publishers, or other indi-\par
viduals or organizations that hold the copyright. To copy an entire book or a part of it, per-\par
mission must be received from the copyright owner, who will most likely expect to be paid.\par
Copyright law distinguishes between different types of intellectual property. Music may\par
be played by anyone after it is published. However, if it is performed for profit, the perform-\par
ers need to pay a fee, called a royalty. A similar principle applies to performances of songs\par
and plays. On the other hand, names, ideas, and book titles are excepted. Ideas do not be-\par
come copyrighted property until they are published in a book, a painting, or a musical\par
work. Almost all artistic work created before the 20th century is not copyrighted because it\par
was created before the copyright law was passed.\par
The two common ways of infringing upon the copyright are plagiarism and piracy. Pla-\par
giarizing the work of another person means passing it off as one's own. The word plagiarism\par
is derived from the Latin plagiarus, which means "abductor" Piracy may be an act of one per-\par
son but, in many cases, it is a joint effort of several people who reproduce copyrighted mate-\par
rial and sell it for profit without paying royalties to the creator. Technological innovations\par
have made piracy easy, and anyone can duplicate a motion picture on videotape, a computer\par
program, or a book. Video cassette recorders can be used by practically anyone to copy mov-\par
ies and television programs, and copying software has become almost as easy as copying a\par
book. Large companies zealously monitor their copyrights for. slogans, advertisements, and\par
brand names, protected by a trademark.\par
\endgroup
\medskip
\noindent\textbf{13.} What does the passage mainly discuss?
\begin{enumerate}[label=(\Alph*),leftmargin=*,itemsep=0.2em]
\item Legal rights of property owners
\item Legal ownership of creative work
\item Examples of copyright piracy
\item Copying creating work for profit
\end{enumerate}
\par
\vspace{0.25\baselineskip}
\noindent\textbf{14.} In lines 3 and 4, the word "extended" is closest in meaning to
\begin{enumerate}[label=(\Alph*),leftmargin=*,itemsep=0.2em]
\item explicated
\item exposed
\item guranteed
\item granted
\end{enumerate}
\par
\vspace{0.25\baselineskip}
\noindent\textbf{15.} It can be inferred from the passage that copyright law is intended to protect
\begin{enumerate}[label=(\Alph*),leftmargin=*,itemsep=0.2em]
\item the user's ability to enjoy an artistic work
\item the creator's ability to profit from the work
\item paintings and photographs from theft
\item computer software and videos from being copied
\end{enumerate}
\par
\vspace{0.25\baselineskip}
\noindent\textbf{16.} In line 16, the word "principle" is closest in meaning to
\begin{enumerate}[label=(\Alph*),leftmargin=*,itemsep=0.2em]
\item crucial point
\item cardinal role
\item fundamental rule
\item formidable force
\end{enumerate}
\par
\vspace{0.25\baselineskip}
\noindent\textbf{17.} Which of the following properties is NOT mentioned as protected by copyright?
\begin{enumerate}[label=(\Alph*),leftmargin=*,itemsep=0.2em]
\item music and plays
\item paintings and maps
\item printed medium
\item scientific discoveries
\end{enumerate}
\par
\vspace{0.25\baselineskip}
\noindent\textbf{18.} It can be inferred from the passage that it is legal if
\begin{enumerate}[label=(\Alph*),leftmargin=*,itemsep=0.2em]
\item two songs, written by two different composers, have the same melody
\item two books, written by two different authors, have the same titles
\item two drawings, created by two different artists, have the same images
\item two plays, created by two different playwrights, have the same plot and characters
\end{enumerate}
\par
\vspace{0.25\baselineskip}
\noindent\textbf{19.} With which of the following statements is the author most likely to agree?
\begin{enumerate}[label=(\Alph*),leftmargin=*,itemsep=0.2em]
\item Teachers are not allowed to make copies of published materials for use by their students.
\item Plays written in the 16th century cannot be performed in theaters without permission.
\item Singers can publicly sing only the songs for which they wrote the music and the lyrics.
\item It is illegal to make photographs when sightseeing or traveling.
\end{enumerate}
\par
\vspace{0.25\baselineskip}
\noindent\textbf{20.} In line 21, the phrase "infringing upon" is closest in meaning to
\begin{enumerate}[label=(\Alph*),leftmargin=*,itemsep=0.2em]
\item impinging upon
\item inducting for
\item violating
\item abhorring
\end{enumerate}
\par
\vspace{0.25\baselineskip}
\noindent\textbf{21.} The purpose of copyright law is most comparable with the purpose of which of the following?
\begin{enumerate}[label=(\Alph*),leftmargin=*,itemsep=0.2em]
\item A law against theft
\item A law against smoking
\item A school policy
\item A household rule
\end{enumerate}
\par
\vspace{0.25\baselineskip}
\noindent\textbf{22.} According to the passage, copyright law is
\begin{enumerate}[label=(\Alph*),leftmargin=*,itemsep=0.2em]
\item meticulously observed
\item routinely ignored
\item frequently debated
\item zealously enforced
\end{enumerate}
\par
\vspace{0.25\baselineskip}
\noindent\textbf{23.} In line 27, the word "practically" is closest in meaning to
\begin{enumerate}[label=(\Alph*),leftmargin=*,itemsep=0.2em]
\item truthfully
\item hardly
\item clearly
\item almost
\end{enumerate}
\par
\vspace{0.25\baselineskip}
\subsection*{Questions 24-33}
\begingroup
\small
\sloppy
\setlength{\parskip}{0pt}
\setlength{\parindent}{0pt}
The official residence of the president of the United States is the White House, located at\par
1600 Pennsylvania Avenue, in Washington, D.C. The Commissioners of the District of Co-\par
lumbia held a meeting in 1792 and decided to hold a contest for the best design for the\par
Presidential House. James Hoban, an architect born in Ireland, was the winner. His bid for\par
the construction of the mansion asked for \$200,000, but the final cost of the building came\par
to twice that amount. The work on the project began during the same year, and the grounds\par
of approximately one and a half miles west of the Capitol Hill were chosen by Major Pierre-\par
Charles L'Enfant, who was in charge of city planning. However, the construction continued\par
for several more years, and George Washington had stepped down as president before the\par
building was habitable. When John Adams, the second president of the United States and\par
his wife Abigail moved in in 1800, only six rooms had been completed.\par
The grey sandstone walls of the house were painted white during construction, and the\par
color of the paint gave the building its name. The building was burned on August 24, 1814,\par
and James Hoban reconstructed the house for President James Monroe and his family, who\par
moved there in 1817. The north portico was added to the building in 1829, water pipes\par
were installed in 1833, gas lighting in 1848, and electricity in 1891. In 1948, inspectors an-\par
nounced that the building was so dilapidated that it was beyond repair and suggested that it\par
was cheaper to construct a new one than repair the existing dwelling. However, the national\par
sentiment was to keep the original form intact, and Congress appropriated \$5.4 million dol-\par
lars for repairs. In 1961, Jacqueline Kennedy launched a program to redecorate the rooms\par
and appointed a Fine Arts Committee to make choices of furnishing and colors.\par
The house of the president accords its residents a great deal of space. The living quarters\par
contain 107 rooms, 40 corridors, and 19 baths. The White House contains a doctor's suite, a\par
dentist's office, a large solarium, a broadcasting room, and a two-floor basement for storage\par
and service rooms. The office in which the president works is not located in the White\par
House, but in a separate building called the West Wing. The White House stands on 16\par
acres of parklike land and overlooks a broad lawn, flower gardens, and wood groves.\par
\endgroup
\medskip
\noindent\textbf{24.} In line 3, the word "contest" is closest in meaning to
\begin{enumerate}[label=(\Alph*),leftmargin=*,itemsep=0.2em]
\item hearir
\item concourse
\item competition
\item computation
\end{enumerate}
\par
\vspace{0.25\baselineskip}
\noindent\textbf{25.} What does the passage imply about the cost of the White House construction?
\begin{enumerate}[label=(\Alph*),leftmargin=*,itemsep=0.2em]
\item It was proposed at the meeting of the commissioners.
\item It did not adhere to the original estimate.
\item It was not included in the architectural design.
\item It was considered excessive for the presidential home.
\end{enumerate}
\par
\vspace{0.25\baselineskip}
\noindent\textbf{26.} In line 6, the word "grounds" is closest in meaning to
\begin{enumerate}[label=(\Alph*),leftmargin=*,itemsep=0.2em]
\item high ground
\item several lots
\item site
\item hills
\end{enumerate}
\par
\vspace{0.25\baselineskip}
\noindent\textbf{27.} It can be inferred from the passage that
\begin{enumerate}[label=(\Alph*),leftmargin=*,itemsep=0.2em]
\item George Washington often used the White House steps
\item George Washington contributed to the White House design
\item George Washington never lived in the White House
\item The White House was excluded from the city planning
\end{enumerate}
\par
\vspace{0.25\baselineskip}
\noindent\textbf{28.} The author of the passage implies that the construction of the main White House building continued
\begin{enumerate}[label=(\Alph*),leftmargin=*,itemsep=0.2em]
\item up to 1800
\item after 1800
\item until 1814
\item until 1792
\end{enumerate}
\par
\vspace{0.25\baselineskip}
\noindent\textbf{29.} In line 17, the word "dilapidated" is closest in meaning to
\begin{enumerate}[label=(\Alph*),leftmargin=*,itemsep=0.2em]
\item ornate
\item run-down
\item old-fashioned
\item obscure
\end{enumerate}
\par
\vspace{0.25\baselineskip}
\noindent\textbf{30.} What can be inferred about the White House from the information in the second paragraph?
\begin{enumerate}[label=(\Alph*),leftmargin=*,itemsep=0.2em]
\item Few changes occurred in the structure in the first half of the 20th century.
\item The building was modernized extensively during one decade.
\item Running water was installed in the second half of the 19th century.
\item Each president added new features to the building's conveniences.
\end{enumerate}
\par
\vspace{0.25\baselineskip}
\noindent\textbf{31.} In line 19, the word "appropriated" is closest in meaning to
\begin{enumerate}[label=(\Alph*),leftmargin=*,itemsep=0.2em]
\item accumulated
\item accosted
\item authorized
\item aggrandized
\end{enumerate}
\par
\vspace{0.25\baselineskip}
\noindent\textbf{32.} In line 20, the word "launched" is closest in meaning to
\begin{enumerate}[label=(\Alph*),leftmargin=*,itemsep=0.2em]
\item lauded
\item lavished
\item began
\item requested
\end{enumerate}
\par
\vspace{0.25\baselineskip}
\noindent\textbf{33.} The passage mentions all of the following White House premises EXCEPT
\begin{enumerate}[label=(\Alph*),leftmargin=*,itemsep=0.2em]
\item hallways
\item kitchen
\item medical offices
\item storage rooms
\end{enumerate}
\par
\vspace{0.25\baselineskip}
\subsection*{Questions 34-42}
\begingroup
\small
\sloppy
\setlength{\parskip}{0pt}
\setlength{\parindent}{0pt}
The organization of the family system in the United States emphasizes monogamy,\par
neolocal residence, an extended family linkage, and a relatively free choice in the selection of\par
mates. Compared to those in many other countries, American families are usually small and\par
isolated. Family and marital roles for men and women overlap considerably, with little dif-\par
ferentiation. Marriage and divorce rates are high, and, according to demographic estimates,\par
95 percent or more of the adult population marry in the course of their lives. The rise in the\par
divorce and remarriage rates has made definitions of family increasingly complex, although\par
several basic patterns have remained stable.\par
The smallest family consists of two persons: a husband and a wife, a parent and a child, or\par
a sister and a brother. Nuclear families include two or more persons who are related by mar-\par
riage, blood, or adoption and who share a common household. A wide range of appropriate\par
kinship patterns have developed in the past two decades, but marriage, an exclusive rela-\par
tionship between two people, is predominant. The only norm that is universally recognized\par
by the society and by law permits a person to have only one spouse in a monogamous family.\par
Serial monogamy has become ubiquitous, and it assumes that persons can have more than\par
one spouse in the course of their lives, as long as there is only one spouse at a time.\par
Social norms pertaining to the choice of marriage partners prescribe appropriate and un-\par
acceptable characteristics. Most societies, including the U.S., grant acceptance to en-\par
dogamous partners who marry within their own racial, ethnic, socioeconomic, educational,\par
age, and religious groups. Traditionally, such norms are justified by common beliefs that\par
those who share similar role expectations, values, attitudes, and customs are less likely to\par
find themselves in conflict. On the other hand, an outcome of endogamous marriages may\par
be that partners' similar socioeconomic levels tend to keep wealth and power within a group\par
of the same class. Similarly, other unique benefits of such unions is that marrying within the\par
same religious grouping may serve to retain the number of its members and, therefore, its\par
political power and influence in the society as a whole.\par
\endgroup
\medskip
\noindent\textbf{34.} According to the passage, a typical family unit in the U.S.
\begin{enumerate}[label=(\Alph*),leftmargin=*,itemsep=0.2em]
\item has a free selection of mates
\item shares a residence with other families
\item lives close to other family members
\item resides apart from its extended family
\end{enumerate}
\par
\vspace{0.25\baselineskip}
\noindent\textbf{35.} In line 4, the word "overlap" is closest in meaning to
\begin{enumerate}[label=(\Alph*),leftmargin=*,itemsep=0.2em]
\item coincide
\item consign
\item collide
\item concur
\end{enumerate}
\par
\vspace{0.25\baselineskip}
\noindent\textbf{36.} The author of the passage implies that American family relationships
\begin{enumerate}[label=(\Alph*),leftmargin=*,itemsep=0.2em]
\item have remained stable
\item have undergone changes
\item have few disparities
\item have little definition
\end{enumerate}
\par
\vspace{0.25\baselineskip}
\noindent\textbf{37.} What is the basis for the definition of a nuclear family in the U.S.?
\begin{enumerate}[label=(\Alph*),leftmargin=*,itemsep=0.2em]
\item Blood relationships o r adoption
\item Two persons
\item Households
\item Marriage
\end{enumerate}
\par
\vspace{0.25\baselineskip}
\noindent\textbf{38.} In line 12, the word "kinship" is closest in meaning to
\begin{enumerate}[label=(\Alph*),leftmargin=*,itemsep=0.2em]
\item ownership
\item sponsorship
\item family connection
\item family services
\end{enumerate}
\par
\vspace{0.25\baselineskip}
\noindent\textbf{39.} The author of the passage implies that the law in the U.S. permits a person to have
\begin{enumerate}[label=(\Alph*),leftmargin=*,itemsep=0.2em]
\item one spouse in the course of a lifetime
\item several spouses in sequential order
\item a series of spouses at one time
\item two spouses during their lifetime
\end{enumerate}
\par
\vspace{0.25\baselineskip}
\noindent\textbf{40.} In line 17, the word "prescribe" is closest in meaning to
\begin{enumerate}[label=(\Alph*),leftmargin=*,itemsep=0.2em]
\item medicate
\item limit
\item determine
\item legalize
\end{enumerate}
\par
\vspace{0.25\baselineskip}
\noindent\textbf{41.} According to the passage, what marriage partners are considered acceptable in most societies?
\begin{enumerate}[label=(\Alph*),leftmargin=*,itemsep=0.2em]
\item Those that have similar social characteristics.
\item Those that have a broad range of life experiences.
\item Those that have a good socioeconomic foundation.
\item Those that have similar attitudes to having children.
\end{enumerate}
\par
\vspace{0.25\baselineskip}
\noindent\textbf{42.} It can be inferred from the passage that endogamous marriages
\begin{enumerate}[label=(\Alph*),leftmargin=*,itemsep=0.2em]
\item help to diversify the society
\item help to preserve the existing social order
\item increase social mobility of individuals
\item decrease the number of conflicting signals.
\end{enumerate}
\par
\vspace{0.25\baselineskip}
\subsection*{Questions 43-50}
\begingroup
\small
\sloppy
\setlength{\parskip}{0pt}
\setlength{\parindent}{0pt}
For many people, mushrooms are strange, colorless, incomprehensible plants that should\par
be avoided. Quaint tales and scary stories surround mushrooms because some are extremely\par
poisonous. In reality, however, mushrooms are fungi that are simple plants without devel-\par
oped roots, leaves, stems, flowers, or seeds. They grow in wetlands, grassy meadows, and\par
woods. Certain types of mushrooms are delicious and are included as ingredients in many\par
recipes and trendy snacks. For example, morels are considered one of the choicest foods, and\par
truffles, related to morels, are highly prized in Europe. Their shape is tubelike, and they re-\par
main entirely underground, a foot or more below the surface in the old days, dogs and pigs\par
were specially trained to hunt them by scent.\par
Mushrooms stand out among other plants because they have no chlorophyll and cannot\par
generate their own nourishment. The part of the fungus that rises above the ground is the\par
fruiting body, and the vegetative part that produces growth is hidden under the ground. It\par
can be usually dug up in the form of dense, white tangled filaments, which, depending on\par
the food supply and moisture, can live for hundreds of years. In fact, mushrooms, as well as\par
the rest of the fungus genus species, are one of the few remaining simple plants that are be-\par
lieved to be among the oldest living organisms. When their environment is not conducive\par
to growth, filaments stop proliferating and can lie dormant for dozens of years.\par
Although mushrooms are rich in flavor and texture, they have little food value. Picking\par
mushrooms requires a thorough knowledge of environments where they are most likely to\par
grow and an ability to tell between edible and poisonous plants. Most mushrooms thrive in\par
temperatures from 68 to 86 (F) with plenty of moisture, and nearly complete darkness\par
produces the best crop. The entire mushroom should be picked, the stem, the cap, and what-\par
ever part that is underground. Brightly colored mushroom caps usually indicate that the\par
plant is not fit for consumption, and the more the mushroom attracts attention, the more\par
poisonous it is. Mushrooms with beautiful red or orange spotted caps that grow under large\par
trees after a good rain are particularly poisonous. If milky or white juices seep from a break\par
in the body of plant, chances are it should not be picked. Old mushrooms with brown caps\par
are also not very safe.\par
\endgroup
\medskip
\noindent\textbf{43.} In line 2, the word "quaint" i s closest in meaning to
\begin{enumerate}[label=(\Alph*),leftmargin=*,itemsep=0.2em]
\item convoluted
\item fanciful.
\item irritating
\item perfunctory
\end{enumerate}
\par
\vspace{0.25\baselineskip}
\noindent\textbf{44.} With which of the following statements is the author of the passage most likely to agree?
\begin{enumerate}[label=(\Alph*),leftmargin=*,itemsep=0.2em]
\item In the old days, when food was scarce, people chose mushrooms as food.
\item Mushrooms should be treated as all other plants.
\item Because they are poisonous, people should stay away from mushrooms.
\item Mushrooms have different forms of roots, stems, and leaves.
\end{enumerate}
\par
\vspace{0.25\baselineskip}
\noindent\textbf{45.} In line 6, the word "trendy" is closest in meaning to
\begin{enumerate}[label=(\Alph*),leftmargin=*,itemsep=0.2em]
\item tender
\item experimental
\item fashionable
\item trusted
\end{enumerate}
\par
\vspace{0.25\baselineskip}
\noindent\textbf{46.} In line 7, the word "Their" refers to
\begin{enumerate}[label=(\Alph*),leftmargin=*,itemsep=0.2em]
\item morels
\item foods
\item truffles
\item morels and truffles
\end{enumerate}
\par
\vspace{0.25\baselineskip}
\noindent\textbf{47.} It can be inferred from the passage that mushrooms multiply mostly by means of
\begin{enumerate}[label=(\Alph*),leftmargin=*,itemsep=0.2em]
\item moisture
\item fruiting bodies
\item nourishment
\item root systems
\end{enumerate}
\par
\vspace{0.25\baselineskip}
\noindent\textbf{48.} The author of the passage implies that mushrooms
\begin{enumerate}[label=(\Alph*),leftmargin=*,itemsep=0.2em]
\item have been known since ancient times
\item are a relatively recent form of plants
\item cannot survive without a good environment
\item have been carefully analyzed
\end{enumerate}
\par
\vspace{0.25\baselineskip}
\noindent\textbf{49.} In line 20, the word "tell" is closest in meaning to
\begin{enumerate}[label=(\Alph*),leftmargin=*,itemsep=0.2em]
\item narrate
\item distinguish
\item say
\item see
\end{enumerate}
\par
\vspace{0.25\baselineskip}
\noindent\textbf{50.} What does the author of the passage imply about brightly colored mushrooms?
\begin{enumerate}[label=(\Alph*),leftmargin=*,itemsep=0.2em]
\item They are beautiful.
\item They should not be eaten.
\item They attract attention.
\item They should be destroyed.
\end{enumerate}
\par
\vspace{0.25\baselineskip}
\end{document}
